\section{Conclusions and Future Work}
\label{sec:conclusion}

The joint between the plasma-facing and structural materials of the fusion first wall
is limited by thermal-expansion mismatch and by the formation of brittle W--Fe
intermetallics. In this work, a compositionally graded alloy was designed to bridge
the two sides within an eleven-element refractory composition space. Activation
screening was performed by linear programming on the composition-linear FLARE model,
CALPHAD calculations were used to restrict the space to single-phase BCC compositions
of adequate thermal conductivity, and the compositions of the resulting design region
were classified as plasma-facing or structural materials by solidus, Pugh ratio, and
EUROFER97-level yield strength. The gradient path between the two terminal alloys was
planned on the composition graph with the lack-of-monotonicity cost function of Kirk
et al., with monotonic CTE and solidus imposed as hard constraints. The primary
conclusions that can be drawn from this work are as follows:

\begin{enumerate}
\item Of the \(61{,}303\) unique activation-feasible compositions, \(700\) are
single-phase BCC and thermally conductive, and \(686\) of these form one connected
design region. Phase stability is by far the most restrictive of the property
criteria.
\item Imposing monotonicity as hard constraints reduces the cost function to path
length, so the optimal path is found exactly by Dijkstra's algorithm, with no
penalty weight to select.
\item The resulting path joins \ch{W90V5Re5} to \ch{W15V75Ti10} in sixteen single
\(5\)~at.\% exchanges, the minimum possible for these terminals, so the monotonicity
constraints are met at no cost in path length.
\item Along the path, the linear CTE at \(600^{\circ}\)C rises from \(5.03\) to
\(9.28\times10^{-6}\)~K\(^{-1}\) in increments of at most
\(0.28\times10^{-6}\)~K\(^{-1}\), the solidus falls from \(3449\) to \(2262\)~K, and
every layer remains single-phase BCC, reduced in activation, and at least as strong as
EUROFER97 at \(600\) and \(650^{\circ}\)C. The Pugh ratio, although not constrained,
also rises at every step, from \(2.00\) to \(3.03\).
\item Despite being planned in an eleven-element space, the path contains only four
elements, a W--V gradient with small Ti and Re additions, which simplifies its
deposition.
\end{enumerate}

The formulation generalizes in both of its parts. The constraint-first screening
applies to any model whose responses are linear in composition, and the path design
applies to any graded alloy that must be monotonic in one or more properties. For
future work in this area, a number of steps could be taken. The CTE, solidus, and
strength of the terminal alloys and selected intermediate layers should be measured
to validate the model estimates. The gradient should then be fabricated by additive
manufacturing and its microstructure characterized along the build. Finally, the
thermal stresses and thermomechanical response of the graded wall under first-wall
heat loads should be evaluated.
% An edge-localized-mode assessment independently validates the design,
% with the plasma-facing terminal carrying a melting margin of \(+1650\)~K and
% the margin remaining positive along the gradient precisely through its
% plasma-facing extent, the structural layers behind it sitting shielded from
% surface loading.
